\documentclass[11pt]{article}
\usepackage[utf8]{inputenc}
\usepackage{amsmath,amssymb}
\usepackage[margin=1in]{geometry}
\usepackage{hyperref}
\emergencystretch=3em

\title{Linearity and the difference form of the canonical coefficients}
\author{Claude, with Daniel Murfet}
\date{2026-09-07}

\begin{document}
\maketitle
\sloppy

\begin{abstract}
Completion of Astra~\#10's R6. On the envelope class the canonical coefficients are linear:
$A_\alpha(c-d)=A_\alpha(c)-A_\alpha(d)$ and $B_\alpha(c-d)=B_\alpha(c)-B_\alpha(d)$
(\texttt{canonA\_sub}, \texttt{canonB\_sub}), proved through the collision formula and the moment
series of unit~126 (the finite-part functionals are linear on absolutely convergent rows; the
envelope class is closed under differences with envelope $H+H'$). With the weighted-norm bounds of
unit~140 this gives the difference form Astra asked for:
$|A_\alpha(c)-A_\alpha(d)|\le w_C(\alpha)\|c-d\|_{\alpha,0,r}$ and
$|B_\alpha(c)-B_\alpha(d)|\le(w_U+w_V)(\alpha)\|c-d\|_{\alpha,0,r}+w_C(\alpha)\|c-d\|_{\alpha,1,r}$
(\texttt{abs\_canonA\_sub\_le}, \texttt{abs\_canonB\_sub\_le}). Zero \texttt{sorry}/\texttt{axiom}.
\end{abstract}

\section{The things formalised}
{\small
\begin{verbatim}
theorem logMoment_sub (β α : ℝ) (ℓ : ℕ) (c₁ c₂ : ℝ → ℝ)
    (h₁ : IntegrableOn ...) (h₂ : ...) :
    logMoment β α ℓ (fun s => c₁ s - c₂ s) = logMoment β α ℓ c₁ - logMoment β α ℓ c₂
theorem sub_envelope (c d : ℕ × ℕ → ℝ → ℝ) (ρ : ℝ) (hρ : 0 < ρ) (H H' : ℝ → ℝ)
    (hc : ∀ ij s, |c ij s| * ρ ^ (ij.1 + ij.2) ≤ H s)
    (hd : ∀ ij s, |d ij s| * ρ ^ (ij.1 + ij.2) ≤ H' s) (ij : ℕ × ℕ) (s : ℝ) :
    |c ij s - d ij s| * ρ ^ (ij.1 + ij.2) ≤ H s + H' s
theorem canonA_sub ... (α : ℝ) (hα : 0 < α) :
    canonA β h₁ h₂ k₁ k₂ (fun i j s => c (i, j) s - d (i, j) s) α
      = canonA β h₁ h₂ k₁ k₂ (fun i j s => c (i, j) s) α
        - canonA β h₁ h₂ k₁ k₂ (fun i j s => d (i, j) s) α
theorem canonB_sub ... (α : ℝ) (hα : 0 < α) :
    canonB β b h₁ h₂ k₁ k₂ (anaFaceU (fun ij s => c ij s - d ij s) b)
        (anaFaceV (fun ij s => c ij s - d ij s) b)
        (fun i j s => c (i, j) s - d (i, j) s) α
      = canonB ... (c) α - canonB ... (d) α
theorem abs_canonA_sub_le ... (α : ℝ) (hα : 0 < α) :
    |canonA ... (c) α - canonA ... (d) α|
      ≤ cWeight r h₁ h₂ k₁ k₂ α * coeffNorm β α 0 r (fun ij s => c ij s - d ij s)
theorem abs_canonB_sub_le ... (α : ℝ) (hα : 0 < α) :
    |canonB ... (c) α - canonB ... (d) α|
      ≤ (uWeight b r h₁ h₂ k₁ k₂ α + vWeight b r h₁ h₂ k₁ k₂ α)
          * coeffNorm β α 0 r (fun ij s => c ij s - d ij s)
        + cWeight r h₁ h₂ k₁ k₂ α * coeffNorm β α 1 r (fun ij s => c ij s - d ij s)
\end{verbatim}
}

\section{Mathematics}

Linearity is immediate for the collision coefficient and for the log moment (Bochner integrals of
integrable integrands). For the face moments the direct definition through the finite part of the
face function is awkward; instead the moment series $k_1^{-1}\sum_jq_j\,M[\alpha;\ell;c_{ij}]$ of
unit~126, valid for $c$, $d$ and $c-d$ (all in the envelope class), reduces linearity to
$M[c_{ij}-d_{ij}]=M[c_{ij}]-M[d_{ij}]$ and the linearity of absolutely convergent sums. The
difference bound is then the boundedness theorem of unit~140 applied to $c-d$.

\section{Paper-facing meaning}

The coefficient functionals $c\mapsto(A_\alpha(c),B_\alpha(c))$ are bounded linear functionals of
the amplitude coefficient array in the weighted moment norm, with explicit constants: the
deterministic half of the continuity argument that turns $\xi_n\Rightarrow G$ into convergence of
the Taylor-tree coefficients. The remaining nonlinear link is the amplitude map
$(x,y)\mapsto c=e^{\beta sx_{00}}(y*E_{x-x_{00}}(\beta s))$, whose local Lipschitz estimate in the
coefficient Banach-algebra norm is Astra~\#10's rank~5.

\section{Lean notes}

\begin{itemize}
\item \texttt{tsum\_sub} is protected: use \texttt{hf.tsum\_sub hg} (dot notation).
\item \texttt{le\_or\_lt} is gone (\texttt{le\_or\_gt}); here the case split was unnecessary once
$\rho>0$ was made a hypothesis.
\item Proving linearity through a series identity valid on a class closed under differences avoids
any linearity lemma for the finite-part functional itself.
\end{itemize}

\end{document}
